\Needspace{20\baselineskip}
\fancyhead[L]{\small COMET OF 1807.}
\begin{center}
COMET OF 1807.
\end{center}

\begin{center}
\fontsize{10}{12}\selectfont\begin{tabular}{c}
\hline
Göttingische gelehrte Anzeigen. Issue 53. Pp. 521..523. 2 April 1808. \\
\hline
\end{tabular}
\end{center}

According to a letter to Prof. \textsc{Gauss} dated 18 March, Inspector \textsc{Bessel} in Lilienthal was fortunate to be able to follow the latest comet until nearly the end of February. We hasten to share these valuable observations here.


\begin{center}\fontsize{10}{12}\selectfont\setlength{\tabcolsep}{3pt}\renewcommand{\arraystretch}{1.15}\begin{tabular}{lrrr}
1808 & Mean time & \shortstack{Apparent\\right ascension} & \shortstack{Apparent\\north declination}\\\hline
Jan. 4 & $9^{\text{u}}\,49^{\text{m}}\,40^{\text{s}}$ & $325^\circ\;17'\;16''1$ & $47^\circ\;13'\;11''1$\\
4 & $11\;34\;32$ & $325\;21\;47.6$ & $47\;13\;40.9$\\
12 & $6\;58\;59$ & $333\;3\;58.2$ & $47\;40\;58.8$\\
21 & $7\;26\;15$ & $341\;19\;13.0$ & $47\;59\;15.3$\\
23 & $6\;49\;6$ & $343\;2\;37.9$ & $48\;2\;5.2$\\
Feb. 19 & $7\;20\;6$ & $3\;43\;35.7$ & $48\;21\;17.1$\\
20 & $7\;23\;15$ & $4\;24\;3.7$ & $48\;21\;33.5$\\
24 & $7\;35\;16$ & $7\;3\;11.1$ & $48\;23\;38.3$\\
\end{tabular}\end{center}

\par\indent Of the declinations observed in February, only the one on the 20th deserves confidence, in Mr. \textsc{Bessel}’s own judgment.

W\SourceLeaf{306}e also add the latest observation by Dr. \textsc{Olbers} in Bremen. This astronomer’s earlier observations are already known from the Monatl. Correspondenz:

\begin{quote}
Feb. 14. $7^{\text{u}}\,30^{\text{m}}\,36^{\text{s}}$ mean time\quad right asc. $0^\circ\;15'\;53''$,\quad north declination $48^\circ\;18'\;10''$
\end{quote}

\par\indent Mr. \textsc{Bessel} used these latest observations to improve his parabolic elements once more and obtained the following results, which differ little from Prof. \textsc{Gauss}’s elements:


\begin{center}\fontsize{9}{11}\selectfont\setlength{\tabcolsep}{3pt}\renewcommand{\arraystretch}{1.15}\begin{tabular}{lr}
\shortstack[l]{Passage through perihelion,\\Paris time 1807 September} .\;.\;.\; & $18{,}82718$\\
Inclination of the orbit .\;.\;.\;.\;.\;.\;.\;.\;.\;.\;.\;.\;.\;.\;.\;.\;.\; & $63^\circ\;14'\;28''1$\\
Ascending node .\;.\;.\;.\;.\;.\;.\;.\;.\;.\;.\;.\;.\;.\;.\;.\; & $266\;36\;51{,}7$\\
Perihelion .\;.\;.\;.\;.\;.\;.\;.\;.\;.\;.\;.\;.\;.\;.\;.\;.\;.\;.\;.\; & $271\;6\;7{,}5$\\
\multicolumn{2}{l}{(both reckoned from the mean equinox)}\\
Logarithm of the least distance .\;.\;.\;.\;.\;.\;.\;.\;.\;.\;.\;.\; & $9{,}8122168$\\
Logarithm of the mean motion .\;.\;.\;.\;.\;.\;.\;.\;.\;.\;.\;.\; & $0{,}2418031$\\
\end{tabular}\end{center}


\Needspace{14\baselineskip}
\par\indent These elements still leave a discrepancy of one minute in the latitudes at the beginning of December. Distributing the error would reduce it considerably. Mr. \textsc{Bessel}, however, preferred to set the parabolic hypothesis aside and fit an elliptical orbit to the observations as closely as possible. Although the nature of the problem precludes a precise determination of the dimensions of this orbit, it nevertheless seems fairly clear from this that the orbit really is elliptical and that the hyperbolic possibility can be excluded. For few comets can even this much be asserted with firm certainty after a single apparition. The elements of the elliptical orbit determined by Mr. \textsc{Bessel} are as follows:


\begin{center}\fontsize{9}{11}\selectfont\setlength{\tabcolsep}{3pt}\renewcommand{\arraystretch}{1.15}\begin{tabular}{lr}
Time of passage through perihelion, September .\;.\;.\;.\;.\;.\; & $18{,}74986$\\
Inclination of the orbit .\;.\;.\;.\;.\;.\;.\;.\;.\;.\;.\;.\;.\;.\;.\; & $63^\circ\;10'\;53''2$\\
Ascending node .\;.\;.\;.\;.\;.\;.\;.\;.\;.\;.\;.\;.\;.\; & $266\;46\;3{,}1$\\
Perihelion .\;.\;.\;.\;.\;.\;.\;.\;.\;.\;.\;.\;.\;.\;.\;.\;.\;.\; & $271\;56\;0{,}1$\\
\multicolumn{2}{l}{(both reckoned from the mean equinox.)}\\
Logarithm of the least distance .\;.\;.\;.\;.\;.\;.\;.\;.\;.\;.\; & $9{,}8105558$\\
Logarithm of the mean motion .\;.\;.\;.\;.\;.\;.\;.\;.\;.\;.\; & $0{,}2442946$\\
Eccentricity .\;.\;.\;.\;.\;.\;.\;.\;.\;.\;.\;.\;.\;.\;.\;.\;.\;.\;.\; & $0{,}9958626$\\
Semimajor axis .\;.\;.\;.\;.\;.\;.\;.\;.\;.\;.\;.\;.\;.\;.\;.\;.\;.\; & $156{,}253$\\
Orbital period .\;.\;.\;.\;.\;.\;.\;.\;.\;.\;.\;.\;.\;.\;.\;.\;.\;.\;.\;.\; & $1953{,}2\text{ years.}$\\
\end{tabular}\end{center}
